\chapter{Introduction}
\label{chap:cz-introduction}


This is a revised version of the article that originally appeared
in Asian J. Math., {\bf 10}({\bf 2}) (2006), 165--492. In July, we
received complaint from Kleiner and Lott for the lack of
attribution of their work caused by our oversight in the prior
version; we have apologized and acknowledged their contribution in
an erratum to appear in December 2006 issue of the Asian Journal
of Mathematics. In this revision, we have also tried to amend
other possible oversights by updating the references and
attributions. In the meantime, we have changed the title and
modified the abstract in order to better reflect our view that the
full credit of proving the Poincare's conjecture goes to Hamilton
and Perelman. We regret all the oversights occurred in the prior
version and hope that this revision will right these inattentions.
Other than these modifications mentioned, the paper remains
unchanged.

In this paper, we shall present the Hamilton-Perelman theory of
Ricci flow. Based on it, we shall give a detailed account of
complete proofs of the Poincar\'e conjecture and the
geometrization conjecture of Thurston. While the results presented
here are based on the accumulated works of many geometric
analysts, the complete proof of the Poincar\'e conjecture is due
to Hamilton and Perelman.

An important problem in differential geometry is to find a
canonical metric on a given manifold. In turn, the existence of a
canonical metric often has profound topological implications. A
good example is the classical uniformization theorem in two
dimensions which, on one hand, provides a complete topological
classification for compact surfaces, and on the other hand shows
that every compact surface has a canonical geometric structure: a
metric of constant curvature.

How to formulate and generalize this two-dimensional result to
three and higher dimensional manifolds has been one of the most
important and challenging topics in modern mathematics. In 1977,
W. Thurston \cite{Th82}, based on ideas about Riemann surfaces,
Haken's work and Mostow's rigidity theorem, etc, formulated a
geometrization conjecture for three-manifolds which, roughly
speaking, states that every compact orientable three-manifold has
a canonical decomposition into pieces, each of which admits a
canonical geometric structure. In particular, Thurston's
conjecture contains, as a special case, the Poincar\'e conjecture:
A closed three-manifold with trivial fundamental group is
necessarily homeomorphic to the 3-sphere $\mathbb{S}^3$. In the
past thirty years, many mathematicians have contributed to the
understanding of this conjecture of Thurston. While Thurston's
theory is based on beautiful combination of techniques from
geometry and topology, there has been a powerful development of
geometric analysis in the past thirty years, lead by S.-T. Yau, R.
Schoen, C. Taubes, K. Uhlenbeck, and S. Donaldson, on the
construction of canonical geometric structures based on nonlinear
PDEs (see, e.g., Yau's survey papers \cite{Y78, Y06}). Such
canonical geometric structures include K\"ahler-Einstein metrics,
constant scalar curvature metrics, and self-dual metrics, among
others. However, the most important contribution for geometric
analysis on three-manifolds is due to Hamilton.

In 1982, Hamilton \cite{Ha82} introduced the Ricci flow
$$\frac{\partial g_{ij}}{\partial t}=-2R_{ij}$$
to study compact three-manifolds with positive Ricci curvature.
The Ricci flow, which evolves a Riemannian metric by its Ricci
curvature, is a natural analogue of the heat equation for metrics.
As a consequence, the curvature tensors evolve by a system of
diffusion equations which tends to distribute the curvature
uniformly over the manifold. Hence, one expects that the initial
metric should be improved and evolve into a canonical metric,
thereby leading to a better understanding of the topology of the
underlying manifold. In the celebrated paper \cite{Ha82}, Hamilton
showed that on a compact three-manifold with an initial metric
having positive Ricci curvature, the Ricci flow converges, after
rescaling to keep constant volume, to a metric of positive
constant sectional curvature, proving the manifold is
diffeomorphic to the three-sphere $\mathbb{S}^3$ or a quotient of
the three-sphere $\mathbb{S}^3$ by a linear group of isometries.
Shortly after, Yau suggested that the Ricci flow should be the
best way to prove the structure theorem for general
three-manifolds. In the past two decades, Hamilton proved many
important and remarkable theorems for the Ricci flow, and laid the
foundation for the program to approach the Poincar\'e conjecture
and Thurston's geometrization conjecture via the Ricci flow.

The basic idea of Hamilton's program can be briefly described as
follows. For any given compact three-manifold, one endows it with
an arbitrary (but can be suitably normalized by scaling) initial
Riemannian metric on the manifold and then studies the behavior of
the solution to the Ricci flow. If the Ricci flow develops
singularities, then one tries to find out the structures of
singularities so that one can perform (geometric) surgery by
cutting off the singularities, and then continue the Ricci flow
after the surgery. If the Ricci flow develops singularities again,
one repeats the process of performing surgery and continuing the
Ricci flow. If one can prove there are only a finite number of
surgeries during any finite time interval and if the long-time
behavior of solutions of the Ricci flow with surgery is well
understood, then one would recognize the topological structure of
the initial manifold.

Thus Hamilton's program, when carried out successfully, will give
a proof of the Poincar\'e conjecture and Thurston's geometrization
conjecture. However, there were obstacles, most notably the
verification of the so called ``Little Loop Lemma" conjectured by
Hamilton \cite{Ha95F} (see also \cite{CCCY}) which is a certain
local injectivity radius estimate, and the verification of the
discreteness of surgery times. In the fall of 2002 and the spring
of 2003, Perelman \cite{P1, P2} brought in fresh new ideas to
figure out important steps to overcome the main obstacles that
remained in the program of Hamilton. (Indeed, in page 3 of
\cite{P1}, Perelman said ``the implementation of Hamilton program
would imply the geometrization conjecture for closed
three-manifolds" and ``In this paper we carry out some details of
Hamilton program".) Perelman's breakthrough on the Ricci flow
excited the entire mathematics community. His work has since been
examined to see whether the proof of the Poincar\'e conjecture and
geometrization program, based on the combination of Hamilton's
fundamental ideas and Perelman's new ideas, holds together. The
present paper grew out of such an effort.


Now we describe the three main parts of Hamilton's program in more
detail.

\vskip 0.1cm \noindent {\bf (i) Determine the structures of
singularities}

Given any compact three-manifold $M$ with an arbitrary Riemannian
metric, one evolves the metric by the Ricci flow. Then, as
Hamilton showed in \cite{Ha82}, the solution $g(t)$ to the Ricci
flow exists for a short time and is unique (also see Theorem
1.2.1). In fact, Hamilton \cite{Ha82} showed that the solution
$g(t)$ will exist on a maximal time interval $[0, T)$, where
either $T=\infty$, or $0<T<\infty$ and the curvature becomes
unbounded as $t$ tends to $T$. We call such a solution $g(t)$ a
maximal solution of the Ricci flow. If $T<\infty$ and the
curvature becomes unbounded as $t$ tends to $T$, we say the
maximal solution develops singularities as $t$ tends to $T$ and
$T$ is the singular time.

In the early 1990s, Hamilton systematically developed methods to
understand the structure of singularities. In \cite{Ha93}, based
on suggestion by Yau, he proved the fundamental Li-Yau \cite{LY}
type differential Harnack estimate (the Li-Yau-Hamilton estimate)
for the Ricci flow with nonnegative curvature operator in all
dimensions. With the help of Shi's interior derivative estimate
\cite{Sh89}, he \cite{Ha95} established a compactness theorem for
smooth solutions to the Ricci flow with uniformly bounded
curvatures and uniformly bounded injectivity radii at the marked
points. By imposing an injectivity radius condition, he rescaled
the solution to show that each singularity is asymptotic to one of
the three types of singularity models \cite{Ha95F}. In \cite
{Ha95F} he discovered (also independently by Ivey \cite{Iv}) an
amazing curvature pinching estimate for the Ricci flow on
three-manifolds. This pinching estimate implies that any
three-dimensional singularity model must have nonnegative
curvature. Thus in dimension three, one only needs to obtain a
complete classification for nonnegatively curved singularity
models.

For Type I singularities in dimension three, Hamilton \cite{Ha95F}
established an isoperimetric ratio estimate to verify the
injectivity radius condition and obtained spherical or necklike
structures for any Type I singularity model. Based on the
Li-Yau-Hamilton estimate, he showed that any Type II singularity
model with nonnegative curvature is either a steady Ricci soliton
with positive sectional curvature or the product of the so called
cigar soliton with the real line \cite{Ha93E}. (Characterization
for nonnegatively curved Type III models was obtained in
\cite{CZ00}.) Furthermore, he developed a dimension reduction
argument to understand the geometry of steady Ricci solitons
\cite{Ha95F}. In the three-dimensional case, he showed that each
steady Ricci soliton with positive curvature has some necklike
structure. Hence Hamilton had basically obtained a canonical
neighborhood structure at points where the curvature is comparable
to the maximal curvature for solutions to the three-dimensional
Ricci flow.

However two obstacles remained: (a) the verification of the
imposed injectivity radius condition in general; and (b) the
possibility of forming a singularity modelled on the product of
the cigar soliton with a real line which could not be removed by
surgery. The recent spectacular work of Perelman \cite{P1} removed
these obstacles by establishing a local injectivity radius
estimate, which is valid for the Ricci flow on compact manifolds
in all dimensions. More precisely, Perelman proved two versions of
``no local collapsing" property (Theorem 3.3.3 and Theorem 3.3.2),
one with an entropy functional he introduced in \cite{P1}, which
is monotone under the Ricci flow, and the other with a space-time
distance function obtained by path integral, analogous to what
Li-Yau did in \cite{LY}, which gives rise to a monotone
volume-type (called reduced volume by Perelman) estimate. By
combining Perelman's no local collapsing theorem I${'}$ (Theorem
3.3.3) with the injectivity radius estimate in Theorem 4.2.2, one
immediately obtains the desired injectivity radius estimate, or
the Little Loop Lemma (Theorem 4.2.4) conjectured by Hamilton.

Furthermore, Perelman \cite{P1} developed a refined rescaling
argument (by considering local limits and weak limits in
Alexandrov spaces) for singularities of the Ricci flow on
three-manifolds to obtain a uniform and global version of the
canonical neighborhood structure theorem. We would like to point
out that our proof of the singularity structure theorem (Theorem
7.1.1) is different from that of Perelman in two aspects: (1) in
Step 2 of the proof, we only prove a weaker version of Perelman's
Claim 2 in section 12.1 of \cite{P1}, namely finite distance
implies finite curvature. Our treatment, with some modifications,
follows the notes of Kleiner-Lott \cite{KL} in June 2003 on
Perelman's first paper \cite{P1}; (2) in Step 4 of the proof, we
give a new approach to extend the limit backward in time to an
ancient solution.

\vskip 0.1cm \noindent {\bf (ii) Geometric surgeries and the
discreteness of surgery times}

After obtaining the canonical neighborhoods (consisting of
spherical, necklike and caplike regions) for the singularities,
one would like to perform geometric surgery and then continue the
Ricci flow. In \cite{Ha97}, Hamilton initiated such a surgery
procedure for the Ricci flow on four-manifolds with positive
isotropic curvature and presented a concrete method for performing
the geometric surgery. His surgery procedures can be roughly
described as follows: cutting the neck-like regions, gluing back
caps, and removing the spherical regions. As will be seen in
Section 7.3 of this paper, Hamilton's geometric surgery method
also works for the Ricci flow on compact orientable
three-manifolds.

Now an important challenge is to prevent surgery times from
accumulating and make sure one performs only a finite number of
surgeries on each finite time interval. The problem is that, when
one performs the surgeries with a given accuracy at each surgery
time, it is possible that the errors may add up to a certain
amount which could cause the surgery times to accumulate. To
prevent this from happening, as time goes on, successive surgeries
must be performed with increasing accuracy. In \cite{P2}, Perelman
introduced some brilliant ideas which allow one to find ``fine"
necks, glue ``fine" caps, and use rescaling to prove that the
surgery times are discrete.

When using the rescaling argument for surgically modified
solutions of the Ricci flow, one encounters the difficulty of how
to apply Hamilton's compactness theorem (Theorem 4.1.5), which
works only for smooth solutions. The idea to overcome this
difficulty consists of two parts. The first part, due to Perelman
\cite{P2}, is to choose the cutoff radius in neck-like regions
small enough to push the surgical regions far away in space. The
second part, due to the authors and Chen-Zhu \cite{CZ05F}, is to
show that the surgically modified solutions are smooth on some
uniform (small) time intervals (on compact subsets) so that
Hamilton's compactness theorem can still be applied. To do so, we
establish three time-extension results (see Assertions 1-3 of the
Step 2 in the proof of Proposition 7.4.1). Perhaps, this second
part is more crucial. Without it, Shi's interior derivative
estimate (Theorem 1.4.2) may not applicable, and hence one cannot
be certain that Hamilton's compactness theorem holds when only
having the uniform $C^{0}$ bound on curvatures. We remark that in
our proof of this second part, as can be seen in the proof of
Proposition 7.4.1, we require a deep comprehension of the
prolongation of the gluing ``fine" caps for which we will use the
recent uniqueness theorem of Bing-Long Chen and the second author
\cite{CZ05U} for solutions of the Ricci flow on noncompact
manifolds.

Once surgeries are known to be discrete in time, one can complete
the classification, started by Schoen-Yau \cite{ScY79m, ScY79},
for compact orientable three-manifolds with positive scalar
curvature. More importantly, for simply connected three-manifolds,
if one can show that solutions to the Ricci flow with surgery
become extinct in finite time, then the Poincar\'e conjecture
would follow. Such a finite extinction time result was proposed by
Perelman \cite{P3}, and a proof also appears in Colding-Minicozzi
\cite{CM}. Thus, the combination of Theorem 7.4.3 (i) and the
finite extinction time result provides a complete proof to the
Poincar\'e conjecture.

\vskip 0.1cm \noindent {\bf (iii) The long-time behavior of
surgically modified solutions.}

To approach the structure theorem for general three-manifolds, one
still needs to analyze the long-time behavior of surgically
modified solutions to the Ricci flow. In \cite{Ha99}, Hamilton
studied the long time behavior of the Ricci flow on compact
three-manifolds for a special class of (smooth) solutions, the so
called nonsingular solutions. These are the solutions that, after
rescaling to keep constant volume, have (uniformly) bounded
curvature for all time. Hamilton \cite{Ha99} proved that any
three-dimensional nonsingular solution either collapses or
subsequently converges to a metric of constant curvature on the
compact manifold or, at large time, admits a thick-thin
decomposition where the thick part consists of a finite number of
hyperbolic pieces and the thin part collapses. Moreover, by
adapting Schoen-Yau's minimal surface arguments in \cite{ScY79}
and using a result of Meeks-Yau \cite{MY}, Hamilton showed that
the boundary of hyperbolic pieces are incompressible tori.
Consequently, when combined with the collapsing results of
Cheeger-Gromov \cite{ChG86, ChG90}, this shows that any
nonsingular solution to the Ricci flow is geometrizable in the
sense of Thurston \cite{Th82}. Even though the nonsingular
assumption seems very restrictive and there are few conditions
known so far which can guarantee a solution to be nonsingular,
nevertheless the ideas and arguments of Hamilton's work
\cite{Ha99} are extremely important.

In \cite{P2}, Perelman modified Hamilton's arguments to analyze
the long-time behavior of arbitrary smooth solutions to the Ricci
flow and solutions with surgery to the Ricci flow in dimension
three. Perelman also argued that the proof of Thurston's
geometrization conjecture could be based on a thick-thin
decomposition, but he could only show the thin part will only have
a (local) lower bound on the sectional curvature. For the thick
part, Perelman \cite{P2} established a crucial elliptic type
estimate, which allowed him to conclude that the thick part
consists of hyperbolic pieces. For the thin part, he announced in
\cite{P2} a new collapsing result which states that if a
three-manifold collapses with (local) lower bound on the sectional
curvature, then it is a graph manifold. Assuming this new
collapsing result, Perelman \cite{P2} claimed that the solutions
to the Ricci flow with surgery have the same long-time behavior as
nonsingular solutions in Hamilton's work, a conclusion which would
imply a proof of Thurston's geometrization conjecture. Although
the proof of this new collapsing result promised by Perelman in
\cite{P2} is still not available in literature, Shioya-Yamaguchi
\cite{ShY} has published a proof of the collapsing result in the
special case when the manifold is closed. In the last section of
this paper (see Theorem 7.7.1), we will provide a proof of
Thurston's geometrization conjecture by only using
Shioya-Yamaguchi's collapsing result. In particular, this gives
another proof of the Poincar\'e conjecture.

We would like to point out that Perelman \cite{P2} did not quite
give an explicit statement of the thick-thin decomposition for
surgical solutions. When we were trying to write down an explicit
statement, we needed to add a restriction on the relation between
the accuracy parameter $\varepsilon$ and the collapsing parameter
$w$. Nevertheless, we are still able to obtain a weaker version of
the thick-thin decomposition (Theorem 7.6.3) that is sufficient to
deduce the geometrization result.

In this paper, we shall give detailed exposition of what we
outlined above, especially of Perelman's work in his second paper
\cite{P2} in which many key ideas of the proofs are sketched or
outlined but complete details of the proofs are often missing.
Since our paper is aimed at both graduate students and researchers
who intend to learn Hamilton's Ricci flow and to understand the
Hamilton-Perelman theory and its application to the geometrization
of three-manifolds, we have made the paper essentially
self-contained so that the proof of the geometrization is readily
accessible to those who are familiar with basics of Riemannian
geometry and elliptic and parabolic partial differential
equations. In particular, our presentation of Hamilton's works in
general follows very closely his original papers, which are
beautifully written. We hope the readers will find this helpful
and convenient and we thank Professor Hamilton for his generosity.

The reader may find many original papers appeared before
Perelman's preprints, particularly those of Hamilton's on the
Ricci flow, in the book ``Collected Papers on Ricci Flow"
\cite{CCCY} in which the editors provided numerous helpful
footnotes. For introductory materials to the Hamilton-Perelman
theory of Ricci flow, we also refer the reader to the recent book
by B. Chow and D. Knopf \cite{CK04} and the forthcoming book by B.
Chow, P. Lu and L. Ni \cite{CLN}. We remark that there have also
appeared several sets of notes on Perelman's work which cover part
of the materials that are needed for the geometrization program,
especially the notes by Kleiner and Lott \cite{KL} from which we
have benefited for the Step 2 of the proof of Theorem 7.1.1.
(After the original version of our paper went to press, Kleiner
and Lott had put up their latest notes \cite{KL2} on Perelman's
papers on the arXiv on May 25, 2006. Also Morgan and Tian posted
their manuscript \cite{MT} on the arXiv on July 25, 2006.) There
also have appeared several survey articles, such as Cao-Chow
\cite{CC99}, Milnor \cite{Mil}, Anderson \cite{An} and Morgan
\cite{Mor}, on the geometrization of three-manifolds via the Ricci
flow.

We are very grateful to Professor S.-T. Yau, who suggested us to
write this paper based on our notes, for introducing us to the
wonderland of the Ricci flow. His vision and strong belief in the
Ricci flow encouraged us to persevere. We also thank him for his
many suggestions and constant encouragement. Without him, it would
be impossible for us to finish this paper. We are enormously
indebted to Professor Richard Hamilton for creating the Ricci flow
and developing the entire program to approach the geometrization
of three-manifolds. His work on the Ricci flow and other geometric
flows has influenced on virtually everyone in the field. The first
author especially would like to thank Professor Hamilton for
teaching him so much about the subject over the past twenty years,
and for his constant encouragement and friendship.

We are indebted to Dr. Bing-Long Chen, who contributed a great
deal in the process of writing this paper. We benefited a lot from
constant discussions with him on the subjects of geometric flows
and geometric analysis. He also contributed many ideas in various
proofs in the paper. We would like to thank Ms. Huiling Gu, a Ph.D
student of the second author, for spending many months of going
through the entire paper and checking the proofs. Without both of
them, it would take much longer time for us to finish this paper.

We also would like to thank Ben Andrews, Simon Brendle, Shu-Cheng
Chang, Ben Chow, Sun-Chin Chu, Panagiota Daskalopoulos, Klaus
Ecker, Gerhard Huisken, Tom Ilmanen, Dan Knopf, Peter Li, Peng Lu,
Lei Ni, Natasa Sesum, Carlo Sinestrari, Luen-fai Tam, Jiaping
Wang, Mu-Tao Wang, and Brian White from whom the authors have
benefited a lot on the subject of geometric flows through
discussions or collaborations.

\enlargethispage{5mm} The first author would like to express his
gratitude to the John Simon Guggenheim Memorial Foundation, the
National Science Foundation (grants DMS-0354621 and DMS-0506084),
and the Outstanding Overseas Young Scholar Fund of Chinese
National Science Foundation for their support for the research in
this paper. He also would like to thank Tsinghua University in
Beijing for its hospitality and support while he was working
there. The second author wishes to thank his wife, Danlin Liu, for
her understanding and support over all these years. The second
author is also indebted to the National Science Foundation of
China for the support in his work on geometric flows, some of
which has been incorporated in this paper. The last part of the
work in this paper was done and the material in Chapter 3, Chapter
6 and Chapter 7 was presented while the second author was visiting
the Harvard Mathematics Department in the fall semester of 2005
and the early spring semester of 2006. He wants to especially
thank Professor Shing-Tung Yau, Professor Cliff Taubes and
Professor Daniel W. Stroock for the enlightening comments and
encouragement during the lectures. Also he gratefully acknowledges
the hospitality and the financial support of Harvard University.

%%%%\mainmatter
